\chapter*{Appendix A: Standard Surgical Patient Positioning}
\addcontentsline{toc}{chapter}{Appendix A: Standard Surgical Patient Positioning}
\markboth{Appendix A: Surgical Positioning}{Appendix A: Surgical Positioning}

Proper patient positioning in the operating room is critical to provide optimal surgical access while ensuring patient safety. The physiological changes associated with anesthesia, combined with the loss of protective reflexes, make patients vulnerable to nerve injuries, pressure alopecia, compartment syndrome, and cardiovascular instability. 

\section*{1. Supine Position (Dorsal Decubitus)}
\textbf{Description:} The patient lies flat on their back, face upward, with arms either tucked at the sides or extended on padded arm boards at an angle of less than 90 degrees.
\begin{itemize}
  \item \textbf{Surgical Access:} Anterior surface of the body, including the head and neck, chest, abdomen, pelvis, and upper/lower extremities.
  \item \textbf{Common Procedures:} Cholecystectomy, appendectomy, thyroidectomy, coronary artery bypass grafting (CABG), herniorrhaphy, and mastectomy.
  \item \textbf{Safety Precautions:}
  \begin{itemize}
    \item Pad the occiput, sacrum, and heels to prevent pressure injuries.
    \item Keep arm extension on boards to $<90^{\circ}$ to prevent brachial plexus stretch injury.
    \item Place a small pillow under the knees to relieve strain on the lumbar spine.
    \item Ensure palms are supinated or facing inward to protect the ulnar nerve at the elbow.
  \end{itemize}
\end{itemize}

\section*{2. Trendelenburg Position}
\textbf{Description:} The patient is in the supine position with the head of the operating table tilted downward (usually $15^{\circ}$ to $30^{\circ}$), placing the head lower than the pelvis.
\begin{itemize}
  \item \textbf{Surgical Access:} Lower abdominal cavity and pelvis. Gravity shifts the viscera cephalad, exposing pelvic organs.
  \item \textbf{Common Procedures:} Hysterectomy, cystectomy, low anterior resection (LAR), and laparoscopic pelvic surgeries.
  \item \textbf{Safety Precautions:}
  \begin{itemize}
    \item Monitor respiratory status closely, as abdominal viscera compress the diaphragm, reducing functional residual capacity (FRC) and compliance.
    \item Monitor for increased intracranial pressure (ICP) and intraocular pressure (IOP).
    \item Use non-slip mattress pads instead of shoulder braces to avoid brachial plexus compression.
  \end{itemize}
\end{itemize}

\section*{3. Reverse Trendelenburg Position}
\textbf{Description:} The patient is in the supine position with the foot of the operating table tilted downward, placing the head higher than the pelvis.
\begin{itemize}
  \item \textbf{Surgical Access:} Upper abdomen, head, and neck. Gravity shifts abdominal viscera caudad, improving exposure of the subdiaphragmatic region.
  \item \textbf{Common Procedures:} Gastrectomy, bariatric surgeries (e.g., gastric bypass, sleeve gastrectomy), fundoplication, and thyroidectomy.
  \item \textbf{Safety Precautions:}
  \begin{itemize}
    \item Secure a padded footboard to prevent the patient from sliding downward.
    \item Monitor for hypotension due to venous pooling in the lower extremities, which decreases venous return and cardiac output.
    \item Apply sequential compression devices (SCDs) to mitigate deep vein thrombosis (DVT) risk.
  \end{itemize}
\end{itemize}

\section*{4. Lithotomy Position}
\textbf{Description:} The patient lies supine with hips and knees flexed, and legs abducted and supported in stirrups or boot-style leg holders.
\begin{itemize}
  \item \textbf{Surgical Access:} Perineum, vagina, urethra, rectum, and anus.
  \item \textbf{Common Procedures:} Transurethral resection of the prostate (TURP), dilation and curettage (D\&C), hemorrhoidectomy, and low anterior resection (LAR).
  \item \textbf{Safety Precautions:}
  \begin{itemize}
    \item Raise and lower both legs simultaneously with two personnel to prevent hip dislocation or sudden hemodynamic shifts.
    \item Avoid pressure on the lateral head of the fibula to prevent compression of the common peroneal nerve, which results in foot drop.
    \item Ensure padding is placed under the sacrum and behind the knees.
    \item Minimize duration in high lithotomy to reduce the risk of lower extremity compartment syndrome.
  \end{itemize}
\end{itemize}

\section*{5. Prone Position}
\textbf{Description:} The patient lies face down on the operating table, supported by chest rolls or a specialized frame to allow free chest excursion.
\begin{itemize}
  \item \textbf{Surgical Access:} Posterior surface of the body, including the spine, buttocks, and posterior aspects of the lower extremities.
  \item \textbf{Common Procedures:} Laminectomy, spinal fusion, pilonidal sinus surgery, and Achilles tendon repair.
  \item \textbf{Safety Precautions:}
  \begin{itemize}
    \item Protect the eyes carefully: ensure eyelids are taped shut and avoid any direct pressure on the globe to prevent central retinal artery occlusion (CRAO).
    \item Maintain cervical spine neutrality; pad and support the face using a specialized prone headrest.
    \item Use chest rolls from the clavicle to the iliac crest to allow diaphragmatic movement and prevent venous engorgement of the spinal epidural plexus.
    \item Pad the knees, elbows, and male genitalia to prevent compression injuries.
  \end{itemize}
\end{itemize}

\section*{6. Lateral Decubitus Position}
\textbf{Description:} The patient lies on their non-operative side (left lateral or right lateral decubitus), with the bottom leg flexed at the hip and knee for stability and the top leg straight.
\begin{itemize}
  \item \textbf{Surgical Access:} Retroperitoneal space, flank, hip, and hemithorax.
  \item \textbf{Common Procedures:} Nephrectomy, lobectomy, thoracoscopy (VATS), hip arthroplasty, and adrenalectomy.
  \item \textbf{Safety Precautions:}
  \begin{itemize}
    \item Place an axillary roll (just distal to the axillary vault) to protect the brachial plexus and prevent compression of the dependent axillary artery.
    \item Place a pillow between the legs to prevent adduction of the upper leg and protect the peroneal nerve of the upper leg and saphenous nerve of the lower leg.
    \item Secure the patient to the table using wide tape or beanbags across the hips.
  \end{itemize}
\end{itemize}

\clearpage
